\noindent\underline{命题 3.3 的证明.}
只需考虑 $X\ne\mathbb P^r$ 的情形. 所断言等价的两个性质都与具体坐标的选择无关, 因而可利用特别选取的坐标描述它们. 由于 $X$ 为 $\mathbb P^r$ 中维数 $n\geq1$ 的光滑子簇, 可选齐次坐标 $X_0,\ldots,X_r$, 使:
\begin{enumerate}
\item[3.3.1.] 点 $x_0\in\mathbb P^r$ 为 $(1,0,\ldots,0)$;
\item[3.3.2.] 对 $i=1,\ldots,r$ 置 $x_i=X_i/X_0$, 从而 $\widehat{\mathcal O}_{\mathbb P^r,x_0}$ 等同于 $k[[x_1,\ldots,x_r]]$, 则在 $\operatorname{Spec}\widehat{\mathcal O}_{\mathbb P^r,x_0}$ 中定义 $X$ 的迹的理想由元素 $g_i$ $(i=1,\ldots,r-n)$ 生成, 它们在完备化中具有形式
\[
g_i=x_{n+i}-h_i(x_1,\ldots,x_n),
\qquad i=1,\ldots,r-n,
\]
其中 $h_i$ 为只含 $x_1,\ldots,x_n$ 的幂级数, 属于 $k[[x_1,\ldots,x_n]]$ 的极大理想的平方;
\item[3.3.3.] 在 (3.1.1) 下, 超平面 $H$ 来自 $g_{r-n}$.
\end{enumerate}

置 $V=\operatorname{Spec}(\widehat{\mathcal O}_{X,x_0})$. 由 3.3.2 有平凡化
\[
(I/I^2)|_V\simeq\mathcal O_V^{r-n}
\]
由
\begin{equation}\tag{3.3.4}
(\Lambda_1,\ldots,\Lambda_{r-n})\longmapsto
\sum_{i=1}^{r-n}\Lambda_i g_i.
\end{equation}
给出. 结合 3.3.3, 得到 $(x_0,H)$ 在 $\mathbb P(N)|_V$ 中的一个邻域 $W$, 以及同构 $V\times\mathbb A^{r-n-1}\simeq W$, 由
\begin{equation}\tag{3.3.5}
\begin{aligned}
V\times\mathbb A^{r-n-1}&\longrightarrow\mathbb P(I/I^2)|_V,\\
(v,(\lambda_1,\ldots,\lambda_{r-n-1}))&\longmapsto
\left(v,\left[g_{r-n}+\sum_{i=1}^{r-n-1}\lambda_i g_i\right]\right),
\end{aligned}
\end{equation}
给出, 其中 $\lambda_i=\Lambda_i/\Lambda_{r-n}$.

特别地, $\mathbb P(N)$ 在 $(x_0,H)$ 处局部环的完备化等同于
\[
k[[x_1,\ldots,x_n,\lambda_1,\ldots,\lambda_{r-n-1}]].
\]
现在在完备化层面用这些坐标表示 $\varphi$:
\begin{equation}\tag{3.3.6}
\varphi(x,\lambda)
=\text{由下列方程定义的超平面 }
\sum_{i=0}^{r}\varphi_i(x,\lambda)T_i=0.
\end{equation}
由于此像超平面包含点 $x$, 而后者的齐次坐标为
\begin{equation}\tag{3.3.7}
x=(1,x_1,\ldots,x_n,h_1(x),\ldots,h_{r-n}(x)),
\end{equation}
故在 $\widehat{\mathcal O}_{\mathbb P(N),(x_0,H)}$ 中有
\begin{equation}\tag{3.3.8}
\varphi_0(x,\lambda)
=-\sum_{i=1}^{n}x_i\varphi_i(x,\lambda)
 -\sum_{j=1}^{r-n}h_j(x)\varphi_{n+j}(x,\lambda)
\end{equation}
此外, 由 (3.1.1), 若
\[
D_i=\frac{\partial}{\partial x_i}\in
\operatorname{Der}_k(
\widehat{\mathcal O}_{\mathbb P^r,x_0},
\widehat{\mathcal O}_{\mathbb P^r,x_0}),
\]
则对 $1\leq i\leq r$ 有
\begin{equation}\tag{3.3.9}
\varphi_i(x,\lambda)
=D_i\!\left(g_{r-n}+
\sum_{j=1}^{r-n-1}\lambda_jg_j\right)(x)
\quad\text{在 }\widehat{\mathcal O}_{\mathbb P(N),(x_0,H)}.
\end{equation}

回顾 $g_i=x_{n+i}-h_i(x_1,\ldots,x_n)$, 得到
\begin{equation}\tag{3.3.10}
\varphi_i(x,\lambda)
=-\frac{\partial}{\partial x_i}
\left(h_{r-n}+\sum_{j=1}^{r-n-1}\lambda_jh_j\right),
\qquad 1\leq i\leq n,
\end{equation}
\begin{equation}\tag{3.3.11}
\varphi_{n+j}(x,\lambda)=\lambda_j,
\qquad 1\leq j\leq r-n-1,
\end{equation}
以及
\begin{equation}\tag{3.3.12}
\varphi_r(x,\lambda)=1.
\end{equation}
\footnote{印本在 (3.3.11) 中给出了不相容的下标范围. (3.3.5) 的坐标、系数公式 (3.3.9) 以及印出的 Jacobian 矩阵中的单位矩阵块决定了此处所列的下标.}

设 $\mathfrak m$ 为 $\widehat{\mathcal O}_{\mathbb P(N),(x_0,H)}$ 的极大理想. 由于对 $i=1,\ldots,r-n$ 有 $h_i\equiv0\pmod{\mathfrak m^2}$, 得到
\begin{equation}\tag{3.3.13}
\varphi_0(x,\lambda)\equiv0\pmod{\mathfrak m^2}
\end{equation}
以及
\begin{equation}\tag{3.3.14}
\varphi_i(x,\lambda)
\equiv-\frac{\partial h_{r-n}}{\partial x_i}
\pmod{\mathfrak m^2},
\qquad i=1,\ldots,n.
\end{equation}